\documentclass[11pt]{article}
\usepackage{fontspec}
\usepackage[french,english]{babel}
\usepackage{amsmath}
\usepackage{amssymb}
\usepackage{amsthm}
\usepackage{mathrsfs}
\usepackage{tikz-cd}
\usepackage{enumitem}

\newcommand{\printedpage}[1]{\marginpar{\footnotesize p.~#1}}
\newcommand{\msp}[1]{\marginpar{\footnotesize [#1]}}
\newcommand{\ilim}{\varinjlim\limits}
\newcommand{\uS}{\underline{S}}
\newcommand{\Ob}{\operatorname{Ob}}
\newcommand{\circlabel}[1]{\tikz[baseline=(c.base)]\node[draw,circle,inner sep=1pt](c){#1};}

\begin{document}

\printedpage{114}
% head: A. Grothendieck : Dérivateurs, Ch. XVIII, Édition des Universités de Montpellier II et Paris VII
% foot: Version 17 décembre 2025

\msp{page 303}\noindent\textbf{Definition 4.13.11} (\footnote[132]{Recast with 4.13.23 etc.}). Let $\mathcal{M}$ be an accessible category, $\pi_0$ a useful cardinal for $\mathcal{M}$, $\pi > \pi_0$ a cardinal. A category $I$ is said to be $\pi_0$-$\pi$-\textit{adapted} if
\begin{enumerate}
\item[a)] $I$ is equivalent to a small category,
\item[b)] $I$ is large relative to $\pi$, and
\item[c)] $I$ is stable under $\varinjlim$ of cardinality $\leqslant \pi$, and large relative to $\pi_0$ (\footnote[133]{\textbf{N.B.} If $\pi = \pi_0$, condition c) is equivalent to saying that $I$ is stable under direct factors.}).
\end{enumerate}
\noindent A direct system $(x_i)_{i\in I}$ of type $I$ in $\mathcal{M}$ is said to be $\pi_0$-$\pi$-\textit{adapted}, if $I$ is $\pi_0$-$\pi$-adapted, i.e. satisfies conditions a) b) c), and if moreover we have
\begin{enumerate}
\item[d)] the functor $i \longmapsto x_i : I \longrightarrow \mathcal{M}$ commutes with the $\varinjlim$ (large relative to $\pi_0$ and of cardinality $\leqslant \pi$) considered in c), and
\item[e)] the $x_i$ are in $\mathcal{M}_\pi$.
\end{enumerate}

\noindent\textbf{Proposition 4.13.12.} \textit{Let $\mathcal{M}$ be a $\pi_0$-accessible category, and $\pi > \pi_0$ a useful cardinal for $\mathcal{M}$. Then every $x \in \Ob\mathcal{M}$ admits a representation as $x \xleftarrow{\ \sim\ } \ilim_{j\in J} x_j$, where the direct system $(x_j)_{j\in J}$ in $\mathcal{M}$ is $\pi_0$-$\pi$-adapted. More precisely, one can take $J = \mathcal{M}_\pi/x$.}

\msp{page 304}\noindent Indeed, the verification of properties a) to d) for the canonical direct system, with $\varinjlim$ isomorphic to $x$, is immediate.

The preceding proposition does not say whether one can take an ordered set for $J$. In this case, condition c) means that for every subset $J'$ of $J$ large relative to $\pi_0$ and of cardinality $\leqslant \pi$, $J$ admits a least upper bound $i_J$ in $J$, and condition d) that we then have $x_{i_J} \xleftarrow{\ \sim\ } \ilim_{i\in J'} x_i$.

The natural idea, to construct a $\pi_0$-$\pi$-adapted direct system $(x_j)_{j\in J}$, with $J$ ordered, starting from the one we have canonically, with $J = \mathcal{M}_\pi/x$, is to take for $I$ the set, ordered by inclusion, of subcategories $J'$ of $J_0$ (small full subcategory of $J$ equivalent to $J$) which are of cardinality $\leqslant \pi$ and large relative to $\pi_0$. (Compare (4.7.3.3), page 165). This ordered set is indeed (small and) \textit{large relative to} $\pi$, if we assume that $\pi^{\pi_0} = \pi$ (cf. proof 4.7.5, p. 168-171).

\msp{page 305}\noindent\textbf{Lemma 4.13.13.} \textit{Let $\pi_0$, $\pi > \pi_0$ be infinite cardinals such that $\pi^{\pi_0} = \pi$, and $J$ a small category, large relative to $\pi$. Let $I$ be the set of subcategories $J'$ of $J$ which are large relative to $\pi_0$, and of cardinality $\leqslant \pi$. Then $I$ is $\pi_0$-$\pi$-adapted.}

It remains to see that for every ordered set $A$ large relative to $\pi_0$ and of cardinality $\leqslant \pi$, if $r \longmapsto J_r$ is an increasing family of elements of $I$ indexed by $A$, then
\[
J \ \overset{\text{def}}{=} \ \bigcup_{r\in A} J_r \ = \ \varinjlim_{r\in A} J_r
\]
is in $I$, i.e. is a subcategory of $I$ of cardinality $\leqslant \pi$ (which is clear) and \textit{large relative to} $\pi_0$. This will follow from the

\printedpage{115}
% head: A. Grothendieck : Dérivateurs, Ch. XVIII, Édition des Universités de Montpellier II et Paris VII
% foot: Version 17 décembre 2025

\noindent\textbf{Lemma 4.13.13.1.} \textit{Let $(J_r)_{r\in A}$ be a direct system of categories large relative to $\pi_0$, indexed by an ordered set $A$ large relative to $\pi_0$. Then $J = \ilim_{r\in A} J_r$ is large relative to $\pi_0$.}

\msp{page 306}\noindent We must see that if $B$ is a set of cardinality $\leqslant \pi_0$, then the diagonal functor
\[
\delta \ : \ J \longrightarrow J^B \tag*{$(*)$}
\]
is cofinal. Now since $A$ is large relative to $\pi_0$, $\operatorname{card} B \leqslant \pi_0$, products indexed by $B$ commute with the $\varinjlim$ of type $A$ (SGA 4 I 9.8 applied to Cat and to $\pi_0$). Thus the arrow $(*)$ is interpreted as the direct limit of the similar arrows
\[
\delta_r \ : \ J_r \longrightarrow J_r^B \ , \tag*{$(**)$}
\]
which are cofinal, hence $\delta$ is cofinal as a direct limit of cofinal functors.

\noindent\textbf{Lemma 4.13.13.2.} \textit{Under the hypotheses of 4.13.13, let $\mathcal{M}$ be a category satisfying $\mathrm{L}_{\pi_0}$, and $(x_j)_{j\in J}$ (where $J \longrightarrow \mathcal{M}$) a direct system in $\mathcal{M}$, indexed by $J$. For every $J' \in I$, put}
\[
x_{J'} \ = \ \varinjlim_{j\in J'} x_j
\]
\msp{page 307}\textit{(which exists in $\mathcal{M}$, since $\mathcal{M}$ satisfies $\mathrm{L}_{\pi_0}$). Then the $(x_{J'})_{J'\in I}$ form a direct system in $\mathcal{M}$, and $I \longrightarrow \mathcal{M}$ commutes with direct limits large relative to $\pi_0$ and of cardinality $\leqslant \pi$. If $\mathcal{M}$ is $\pi_0$-accessible, and the $x_j$ ($j \in \Ob J$) are in $\mathcal{M}_\pi$, then $(x_{J'})_{J'\in I}$ is a $\pi_0$-$\pi$-adapted direct system in $\mathcal{M}$.}

All this is essentially trivial, and gives:

\noindent\textbf{Corollary 4.13.14} (of 4.13.12). \textit{Under the conditions of 4.13.12, if $\pi^{\pi_0} = \pi$, one can for every $x \in \Ob\mathcal{M}$ find a representation $x \simeq \ilim_{i\in I} x_i$, where $(x_i)_{i\in I}$ is a $\pi_0$-$\pi$-adapted direct system with $I$} an ordered set. \textit{Moreover, if $\pi' > \pi$ and if $x \in \mathcal{M}_{\pi'}$, one can take $I$ such that we have}
\[
\operatorname{card} I \ \leqslant \ {\pi'}^{\pi} \ .
\]

This last addendum is immediate,\msp{page 308} via the usual estimates.

\noindent\textbf{Theorem 4.13.15} (\footnote[134]{Cf. more general and more symmetric version of theorem 4.13.15 in 4.13.22.}). \textit{Let $\mathcal{M}$ be a $\pi_0$-accessible category, and $\pi > \pi'_0 \geqslant \pi_0$ such that $\pi^{\pi_0} = \pi$} (\footnote[135]{One can restrict oneself to assuming $\pi$ useful for $\mathcal{M}$ (instead of $\pi^{\pi_0} = \pi$), then one finds that a, a$'$, b, c are equivalent and imply c$'$, but are perhaps not equivalent to c$'$.})\textit{. Let $\uS$ be a strictly full subcategory of $\mathcal{M}$, and let $\uS(\pi) = \uS \cap \mathcal{M}_\pi$. The following conditions are equivalent.}
\begin{enumerate}
\item[\circlabel{a}] \textit{$\uS$ is $\pi$-accessible in $\mathcal{M}$ (4.13.3) and $\uS$ is stable in $\mathcal{M}$ under direct limits large relative to $\pi'_0$.}
\end{enumerate}

\printedpage{116}
% head: A. Grothendieck : Dérivateurs, Ch. XVIII, Édition des Universités de Montpellier II et Paris VII
% foot: Version 17 décembre 2025

\begin{enumerate}
\item[a$'$)] \textit{$\uS$ is $\pi$-accessible in $\mathcal{M}$, and $\uS(\pi)$ is stable in $\mathcal{M}_\pi$ under direct limits large relative to $\pi'_0$ and of cardinality $\leqslant \pi$} (\footnote[136]{give a name to this last condition, say $\mathrm{L}(\pi'_0, \pi)$ (for $\uS(\pi)$ in $\mathcal{M}_\pi$).})\textit{.}
\item[b)] \textit{For every object $x$ of $\mathcal{M}$, we have this: the strictly full subcategory $\uS(\pi)/x$ of $J = \mathcal{M}_\pi/x$ is cofinal in $\mathcal{M}_\pi/x$ if and only if $x \in \Ob\uS$. Moreover, $\uS(\pi)$ is stable in $\mathcal{M}_\pi$ under $\varinjlim$ large relative to $\pi'_0$ and of cardinality $\leqslant \pi$.}
\item[\circlabel{c}] \textit{For every object $x$ of $\mathcal{M}$ and for every representation $x \simeq \ilim_{i\in I} x_i$, where $(x_i)_{i\in I}$ is a $\pi'_0$-$\pi$-adapted direct system, (4.13.11), we have this: the strictly full subcategory $I_{\uS}$ of $I$ formed by the $i \in \Ob I$ such that $x_i \in \Ob\uS$, is cofinal in $I$ if and only if $x \in \Ob\uS$. Moreover, $S(\pi)$ is stable in $\mathcal{M}_\pi$ under $\varinjlim$ large relative to $\pi'_0$ and of cardinality $\leqslant \pi$.} % [sic] S(\pi) -> \underline{S}(\pi), printed without the underline used everywhere else for this subcategory
\end{enumerate}
\msp{page 309}
\begin{enumerate}
\item[c$'$)] \textit{As in d)} \texttt{[rather c)]}\textit{, with $I$ an ordered set.}
\end{enumerate}
\noindent\textbf{Corollary 4.13.16.} \textit{The hypotheses on $\mathcal{M}$, $\pi_0$, $\uS$ being those of 4.13.15, for there to exist cardinals $\pi > \pi'_0 \geqslant \pi_0$ satisfying the conditions $\pi^{\pi'_0} = \pi$ (a fortiori $\pi^{\pi_0} = \pi$) and the equivalent conditions a) to d$'$)} \texttt{[rather c$'$)]} \textit{of 4.13.15, it (is necessary and it) suffices that $\uS$ be accessible in $\mathcal{M}$. More precisely, if $\uS$ is accessible in $\mathcal{M}$, let $\pi'_0 \geqslant \pi_0$ be a cardinal such that $\uS$ is $\pi'_0$-accessible in $\mathcal{M}$, and let $\pi$ be a cardinal $> \pi'_0$ such that $\pi^{\pi'_0} = \pi$. Then $(\uS, \pi'_0, \pi)$ satisfies the equivalent conditions of 4.13.15.}

Indeed, the ``it is necessary'' is obvious from condition a) of 4.13.15. For the ``it suffices'' we note that $\uS$ $\pi'_0$-accessible in $\mathcal{M}$ implies that $\uS$ is stable in $\mathcal{M}$ under $\varinjlim$ large relative to $\pi'_0$. Thus to verify a), it suffices to verify that $\uS$ is $\pi$-accessible in $\mathcal{M}$. Now since $\uS$ is\msp{page 310} $\pi'_0$-accessible and $\pi > \pi'_0$ satisfies $\pi^{\pi'_0} = \pi$, we conclude from this that $\uS$ is $\uS_\pi$-accessible, i.e. $\uS \xleftarrow{\ \sim\ } \mathrm{Ind}_\pi(\uS_\pi)$. Thus by 4.13.3, it remains to see that $\uS_\pi \subset \mathcal{M}_\pi$, or what amounts to the same thing, that $\uS_\pi \subset \uS(\pi) \overset{\text{def}}{=} \uS \cap \mathcal{M}_\pi$. (In fact, we shall even have equality, for it is clear that $\uS(\pi) \subset \uS_\pi \ldots$). Now $\uS(\pi)$ is a subcategory of $\uS$ stable under $\varinjlim$ large relative to $\pi'_0$ and of cardinality $\leqslant \pi$ (since $\uS$ and $\mathcal{M}_\pi$ are stable under these limits \textit{in} $\mathcal{M}$, and these limits in $\uS$ and in $\mathcal{M}$ are likewise), and $\uS(\pi)$ contains $S_{\pi'_0}$. Thus $\uS(\pi)$ contains $\uS_\pi$, by 4.7.1 (page 161), implication c) $\Longrightarrow$ b$'$), applied to $\uS$, $\pi'_0$, $\pi$ instead of $\mathcal{M}$, $\pi_0$, $\pi$. This therefore proves corollary 4.13.16. % [sic] S_{\pi'_0} -> \underline{S}_{\pi'_0}, printed without the underline

\textsc{Proof of} 4.13.15. Here are the $\pm$ tautological implications,
\[
\begin{tikzcd}[row sep=small]
 & & & \mathrm{c}' \\
\mathrm{a} \arrow[r, Rightarrow] & \mathrm{a}' \arrow[r, Rightarrow, "\bullet"] & \mathrm{c} \arrow[ur, Rightarrow] \arrow[dr, Rightarrow] & \\
 & & & \mathrm{b}\ ,
\end{tikzcd}
\tag*{$(*)$}
\]
\msp{page 311}which I am going to spell out a little. We have a $\Longrightarrow$ a$'$ since if a is satisfied, $\uS$ \textit{and} $\mathcal{M}_\pi$ are stable \textit{in} $\mathcal{M}$ under $\varinjlim$ large relative to $\pi'_0$ and of cardinality $\leqslant \pi$, hence also their intersection $\uS(\pi)$,
\printedpage{117}
% head: A. Grothendieck : Dérivateurs, Ch. XVIII, Édition des Universités de Montpellier II et Paris VII
% foot: Version 17 décembre 2025
which is therefore also stable in $\mathcal{M}_\pi$ under the said limits. The implication c $\Longrightarrow$ c$'$ is a tautology, and c $\Longrightarrow$ b as well, in view of 4.13.12, which tells us that the canonical direct system $(x_i)_{i\in I}$, with $I = \mathcal{M}_\pi/x$, with limit $x$ is $\pi_0$-$\pi$-adapted. (\textbf{N.B.} Since $\pi^{\pi'_0} = \pi$ and $\pi'_0 \geqslant \pi_0$, we also have $\pi^{\pi_0} = \pi$, hence $\pi$ is useful for $\mathcal{M}$, so that 4.13.12 applies.)

It remains to prove the ``$\pm$ tautology'' a$'$ $\Longrightarrow$ c. So let us take a $\pi_0$-$\pi$-adapted representation
\[
x \ \simeq \ \varinjlim_{i\in I} x_i \ ,
\]
let us prove that
\[
x \ \in \ \Ob\uS \qquad\quad \Longleftrightarrow \qquad\quad I_{\uS} \text{ cofinal in } I.
\]
The implication $\Longleftarrow$ comes from the fact that, $\uS$ being $\pi$-accessible in $\mathcal{M}$, is stable in $\mathcal{M}$ under $\varinjlim$ large relative to $\pi$ (\footnote[137]{We use here, for $\Longleftarrow$, of the hypothesis $(x_i)_{i\in I}$ $\pi_0$-$\pi$-adapted, only the sole fact that $I$ is large relative to $\pi$.}). It remains to prove\msp{page 312} that if $x \in \Ob\uS$, then $I_{\uS}$ is cofinal in $I$. So let us consider the diagram of categories
\[
\begin{tikzcd}
I \arrow[r, "\varphi"] & J = \mathcal{M}_\pi/x \\
I' = I_{\uS} \arrow[u, hook] \arrow[r, "\varphi'"] & J' = \uS(\pi)/x \ , \arrow[u, hook, "\alpha"']
\end{tikzcd}
\]
this is a cartesian square ($I' = I_{\uS}$ is the inverse image under $\varphi : I \longrightarrow J$ of the strictly full subcategory $J'$ of $J$). We have the following properties
\begin{enumerate}
\item[1)] $I$ is large relative to $\pi$, and $\varphi$ cofinal ($\Longrightarrow$ $J$ large relative to $\pi$).
\item[2)] $I$ stable under $\varinjlim$ large relative to $\pi'_0$ and of cardinality $\leqslant \pi$, and $\varphi$ commutes with these $\varinjlim$ .
\item[3)] $J'$ is stable under $\varinjlim$ of the preceding type and $\alpha$ commutes with them (\footnote[138]{This is where the stability hypothesis made in b) comes in, for by b), $\uS$ is $\pi$-accessible in $\mathcal{M}$.}).
\item[4)] $J'$ cofinal in $J$.
\end{enumerate}
\noindent We have just concluded that $I' = \varphi^{-1}(J')$\msp{page 313} is cofinal in $I$. As it is a full subcategory (and even strictly full), it suffices to prove property F 1, i.e. that for every $i \in \Ob I$, exists $i' \in \Ob I'$ and an arrow $i \longrightarrow i'$.

We are going to construct $i'$ again by transfinite construction, on the ordinal interval $I_{\alpha_\pi}$, by constructing a direct system
\[
(i_\alpha)_{\alpha\in I_{\alpha_\pi}} \ , \qquad\quad i_\alpha \in \Ob I \ ,
\]
of objects of $I$, and a direct system $(j_\alpha)_{\alpha\in I_{\alpha_\pi}}$, $j_\alpha \in \Ob J'$, transition arrows denoted $\tau_{\beta,\alpha} : i_\alpha \longrightarrow i_\beta$, $\tau'_{\beta,\alpha} : j_\alpha \longrightarrow j_\beta$, finally for every $\alpha$ in $I_{\alpha_\pi}$ arrows
\[
\begin{tikzcd}
\varphi(i_\alpha) \arrow[r, "u_\alpha"] \arrow[rr, dotted, bend right, "\varphi(\tau_{\alpha+1,\alpha})"'] & j_\alpha \arrow[r, "v_\alpha"] \arrow[rr, dotted, bend left, "\tau'_{\alpha+1,\alpha}"] & \varphi(i_{\alpha+1})\ ( \arrow[r, dotted, "u_{\alpha+1}"] & j_{\alpha+1}\ )
\end{tikzcd}
\tag*{$(*)$}
\]
\printedpage{118}
% head: A. Grothendieck : Dérivateurs, Ch. XVIII, Édition des Universités de Montpellier II et Paris VII
% foot: Version 17 décembre 2025
in $J$, such that
\[
\left\{
\begin{aligned}
v_\alpha u_\alpha \ &= \ \varphi(\tau_{\alpha+1,\alpha}) \\
u_{\alpha+1} v_\alpha \ &= \ \tau'_{\alpha+1,\alpha}
\end{aligned}
\right.
\qquad\qquad \text{(\footnotemark[139])} \ .
\tag*{$(**)$}
\]
\footnotetext[139]{cf. also $(****)$, next page.}

\msp{page 314}\noindent Let us show how the construction of such a system of objects and arrows with
\[
i_0 \ = \ i \tag*{$(***)$}
\]
will imply the existence of $i' \in \Ob I'$ and of $i \longrightarrow i'$. Indeed, since $I_{\alpha_\pi}$ is large relative to $\pi'_0$ (it is large relative to every cardinal $< \pi$), there exists
\[
i' \ \overset{\text{def}}{=} \ \varinjlim_{i\in I_{\alpha_\pi}} i_\alpha % [sic] i\in I_{\alpha_\pi} -> \alpha\in I_{\alpha_\pi}, the index of the limit is \alpha
\]
in $I$, and we have a transition arrow
\[
i \ = \ i_0 \longrightarrow i' \ ,
\]
so it remains to prove that $i' \in \Ob I'$, i.e. $\varphi(i') \in \Ob J'$. Now $\varphi$ commuting with the $\varinjlim$ of type $I_{\alpha_\pi}$, we shall have
\[
\varphi(i) \ = \ \varinjlim \varphi(i_\alpha) \ .
\]
I still need compatibilities
\[
\begin{tikzcd}
\varphi(i_\alpha) \arrow[r, "u_\alpha"] \arrow[d, "\varphi(\tau_{\beta,\alpha})"'] & j_\alpha \arrow[d, "\tau'_{\beta,\alpha}"] & {} \arrow[d, phantom, "\text{and}" description] & j_\alpha \arrow[r, "v_\alpha"] \arrow[d, "\tau'_{\beta,\alpha}"'] & \varphi(i_{\alpha+1}) \arrow[d, "\varphi(\tau_{\beta+1,\alpha+1})"] \\
\varphi(i_\beta) \arrow[r, "u_\beta"] & j_\alpha & {} & j_\beta \arrow[r, "v_\beta"] & \varphi(i_{\beta+1})
\end{tikzcd}
\tag*{$(****)$} % [sic] j_\alpha (lower right corner of the first square) -> j_\beta, target of \tau'_{\beta,\alpha} and of u_\beta
\]
forgotten earlier.

\msp{page 315}\noindent The compatibilities $(****)$ imply that we have homomorphisms
\[
\begin{tikzcd}
\varphi(i) \arrow[r, shift left, "u"] & j \arrow[l, shift left, "v"]
\end{tikzcd}
\qquad\qquad \text{where } j \overset{\text{def}}{=} \varinjlim_{\alpha\in I_{\alpha_\pi}} j_\alpha \quad \text{(limit \textit{in} $J'$, hence in $J$)} \ ,
\]
and the compatibilities $(**)$ imply that
\[
vu \ = \ \mathrm{id}_{\varphi(i)} \ , \qquad\quad uv = \mathrm{id}_J \ , % [sic] \mathrm{id}_J -> \mathrm{id}_j, uv is an endomorphism of j
\]
hence that $u$, $v$ are isomorphisms inverse to each other. Since $J'$ is a \textit{strictly} full subcategory of $J$, it does indeed follow that $\varphi(i) \in J'$, i.e. $i \in I'$, q.e.d.

So everything comes down to constructing the
\[
\underbrace{i_\alpha}_{\in \Ob I} \ , \ \tau_{\beta,\alpha} \ , \qquad\quad \underbrace{j_\alpha}_{\in \Ob J'} \ , \ \tau'_{\beta,\alpha} \ , \qquad\quad u_\alpha \ , \ v_\alpha \ ,
\]

\end{document}
